\chapter{Execution paths}\label{ch:executing}

The same native program runs on two paths (\cref{fig:bypass}): through Metal, which builds the pipeline and submits
the work, and below Metal, where the project's runtime does both and submits through IOGPU
(\cref{tab:native-path-ownership} lists what remains Apple's).

\section{Execution through Metal}\label{ch:through-metal}

The executor, \code{tools/g17decodegen.m}, loads each bundle's library and
binary archive. It builds the pipeline with \code{MTLPipelineOptionFailOnBinaryArchiveMiss}, so Metal must use the
archived code. The GPU then executes this project's code: patching one constant in GPU memory after the pipeline is built
changes the output as predicted (measured, MM~\mm{25.147}), though whether the bytes are copied or checked again
before execution is not determined. Metal still derives the pipeline configuration from the
metadata the image declares and submits the work; the native path (\cref{ch:below-metal-native}) does both itself.

decodegen encodes each token into one command buffer. It creates one Metal buffer per arena, writes the arenas'
initial contents, and encodes the whole token, every dispatch of every layer, into that command buffer. Stage outputs stay on the
GPU, in arenas that the next stage binds. In pipelined mode the GPU also runs the generation step, which advances
the position and writes the next token's embedding. The host then commits the token command buffers back to back
with no wait between them, and reads the generated tokens from a log at the end. With no idle gap between tokens the
GPU stays at its maximum clock state (MM~\mm{0.15}).

In the current InternLM2.5-1.8B decode graph each of the 24 layers is seven dispatches:

\begin{lstlisting}
RMSNorm -> qkv projection -> attention (RoPE + KV append + 32 key slices + merge) -> wo projection + residual
        -> RMSNorm -> w1/w3 projections + SwiGLU -> w2 projection + residual
\end{lstlisting}

Four more dispatches run per token: the final norm, the vocabulary head, the argmax and the generation step. That
makes 172 dispatches a token (MM~\mm{25.154}, MM~\mm{25.190}). mlx-lm runs 16 dispatches a layer for the same model
(MM~\mm{25.142.1}). \emph{Correction:} an earlier graph had eight dispatches a layer (MM~\mm{25.142.1}); MM~\mm{25.211} repeated that figure
and is corrected in place.

Decode speed comes from these design choices:

\begin{itemize}
\item Split-K quantized matrix-vector products read the 4-bit weights with K split across lanes and
  threadgroups, and fuse the residual add or SwiGLU into the same dispatch (MM~\mm{25.152}, MM~\mm{25.153}).
\item The attention is sliced, with a parallel merge. It splits the keys into 32 slices, each a simdgroup with
  its own running softmax, and merges the partial results with a parallel butterfly. Its cost grows with the keys
  it reads, not with a serial merge: 9.7~µs a layer at 128 keys, against MLX's 12.0 (MM~\mm{25.140}, MM~\mm{25.141.14}).
\item The structure is static. Every buffer, binding and dispatch is fixed before the first token, so a token costs one
  command-buffer submission.
\end{itemize}

Metal twins of these kernels, with identical threads, partitions and operation order and compiled by Apple, are
bit-identical and 10–15 percent faster (\cref{sec:bounds-real-workloads}), so the speed comes from these designs and
this compiler's code generation loses part of it. Which code difference costs the time is open (MM~\mm{25.211}).

Batched decode uses the tensor unit. With up to 16 sequences, the batched projection dequantizes each
4-bit weight once in registers, then multiplies it against every sequence with tensor MMAs. It computes the
transpose, \code{Y\textasciicircum{}T = W\textasciicircum{}T X\textasciicircum{}T}, because only B can advance through the key-block index register (MM~\mm{25.166}).

\begin{itemize}
\item The dequantization runs in fp16, on the register-form fp16 fma (\opn{798}) (MM~\mm{25.196}).
\item Splitting K to about 512 threadgroups keeps the GPU occupied (MM~\mm{25.185}).
\item A 32-sequence form runs two 16-row halves over one dequantized weight (MM~\mm{25.178}).
\end{itemize}

Prefill runs the prompt in chunks of up to 512 rows. It uses tensor-unit GEMMs on persistent fp16 weights, an
attention that keeps its score and probability tiles in registers, and an fp16 FFN intermediate. Short prompts use a
16-row 4-bit route instead (MM~\mm{25.163}, MM~\mm{25.183}, MM~\mm{25.202}).

\section{Native execution below Metal}\label{ch:below-metal-native}

The native runtime runs a model with no Metal in the process: it never creates an \code{MTLDevice} or a Metal command
buffer, and it checks that AGXMetal, Apple's Metal driver library, is absent. Work is submitted through Apple's private
\code{IOGPU} framework to the installed kernel driver and the GPU firmware. Two pinned models share one pipeline:
MiniLM-L6 for sentence embeddings, and Qwen2.5-0.5B-Instruct for text generation (MM~\mm{25.210}).

\Cref{tab:native-path-ownership} sets out what is the project's and what is still Apple's (MM~\mm{25.210.9}).

\begin{xltabular}{\linewidth}{@{}lXX@{}}
\caption{What the project supplies at each layer of the native path, and what is still Apple's or outside the project.}\label{tab:native-path-ownership}\\
\toprule
Layer & This project & Still Apple's, or outside \\
\midrule
\endfirsthead
\tablecontinued{3}
\toprule
Layer & This project & Still Apple's, or outside \\
\midrule
\endhead
model & graph import, checkpoint conversion, lowering, memory plan & the checkpoints and tokenizers; NumPy and PyTorch for host input and
validation \\*
GPU program & every instruction byte, from the ordinary compiler & encoding rules recovered with Apple's decoder \\*
execution & the host worker: uploads, resource and command state, Submit, wait,
checks & \code{IOGPU.framework}, IOKit, the AGX kernel driver, the firmware, the
hardware \\*
\bottomrule
\end{xltabular}

The pipeline has five stages (MM~\mm{25.210.3}):

\begin{enumerate}
\def\labelenumi{\arabic{enumi}.}
\item Import: \code{agxforge/g17/inferencegraph.py} reads the checkpoint's tensors into a typed graph, refusing holes,
  overlaps and mismatched shapes.
\item Lower: \code{tools/g17inferencelower.py} compiles every node through the ordinary compiler. MiniLM becomes 170
  stages from 16 distinct bodies. Qwen becomes 748 prefill and 916 decode stages from 28.
\item Plan: a memory planner assigns arenas, each slot with 256-byte guards.
\item Bundle: the schedule is placed into a measured resource graph: 128 MiB for MiniLM, 1 GiB for Qwen.
\item Execute: a resident worker runs it, with at most three bindings per stage and one Submit per stage.
\end{enumerate}

For each stage the runner:

\begin{enumerate}
\def\labelenumi{\arabic{enumi}.}
\item places the stage's program in allocation 0;
\item writes the stage's launch words: the program counter, the shader state, the grid, the threads per threadgroup and
  up to three buffer pointers;
\item builds one count-one Submit record with fresh callback blocks;
\item calls IOGPU and waits, for at most five seconds, for both callbacks.
\end{enumerate}

\Cref{ch:submission-below-metal} documents each of those fields and records, what is known about them, and where the measured envelope
ends. The runtime refuses anything outside that envelope.

Completion and correctness are checked separately.

\begin{itemize}
\item Each stage reply carries preservation bits: the stage completed; the weights, guards, code and uploaded source are
  unchanged; and AGXMetal is absent.
\item Every live region's guards are checked after every stage, and any change in the GPU's recovery or event counters
  fails the run.
\item The captured outputs are then checked offline, at two levels. Each level's bound was registered before the first
  dispatch and never widened (MM~\mm{25.210.2}).

  \begin{enumerate}
  \def\labelenumi{\arabic{enumi}.}
  \item Each stage is checked against its actual input: it is recomputed by independent CPU arithmetic from the
  values its GPU predecessor actually produced. The bound is \code{2e-5 * (1 + abs(reference))}, and packs, gathers,
  bias and residual adds must be bit-exact. All 680 MiniLM stages and 9,984 Qwen stage checks pass. This level checks
  each stage's arithmetic but not drift across the model.
  \item Each application output is checked against the original checkpoint evaluated in FP64.

  \begin{itemize}
  \item MiniLM's final embedding must be within a maximum error of 0.005 and a cosine of at least 0.999. Measured:
  1.29e-4, cosine at least 0.9999997786. Against an FP64 model of the same half-transport policy, the bound
  \code{2e-4 * (1 + abs(reference))} also passes (maximum 8.21e-5).
  \item Each Qwen last-row logit vector, with the native tokens teacher-forced, must be within
  \code{0.05 + 0.003 * abs(reference)}. Measured: 37 of 37 inside, and 37 of 37 agree on the top token.
  \end{itemize}
  \end{enumerate}

  \begin{itemize}
  \item One auxiliary control fails and is reported as failing: Qwen's prefix logits against an FP64 model of the same
  half-projection policy, \code{0.025 + 0.002 * abs(reference)}. The maximum is 0.097, with 10,658 elements outside,
  traced to a half-rounding boundary in layer 0.
  \end{itemize}
\item The numerical domain is half-transported matrix operands with fp32 accumulation, not original-fp32 arithmetic.
\end{itemize}

The native path is slower (measured, MM~\mm{25.210.5}):

\begin{itemize}
\item The matched MiniLM blocks take 2.427 against 1.797~ms for the FFN, and 3.815 against 1.962~ms for attention.
\item In the paired comparison with Transformers on MPS, native Qwen decode ran 3.26–3.33 tok/s against MPS's 22.6–49.1;
  the lowest MPS figure is its first request, which includes first-use effects, and its warm requests ran 48–49.
  Later native changes raised decode to 5.15–5.32 tok/s, paired against the old guard check, and 4.62 tok/s in an
  unpaired replay. Those were not paired with MPS, so the two sets of numbers are not one comparison
  (MM~\mm{25.210.5}).
\item A repository-authored full-model Metal control runs the same 28 programs. Its own decode frame is 230~ms, of which
  the reported GPU time is 33~ms, so the control's frame is dominated by time outside its GPU interval; this does
  not apportion the native runtime's time.
\item Some native costs are measured directly:

  \begin{itemize}
  \item In the FFN block, a warm request spends 2.385~ms from Submit to the completed callback. That interval mixes GPU
  execution, driver work and completion delivery, and is not GPU time. Only 0.0495~ms then passes from callback to
  the waiter's return, so semaphore wakeup is not the cost.
  \item Replacing the host's scalar guard check with a vectorised one made Qwen decode 1.30–1.31× faster.
  \item The rest of the per-token breakdown, across 916 Submits, guard scans and padded transfers, is not measured.
  \end{itemize}
\end{itemize}

\section{Comparison of the two paths}\label{ch:comparing-two-paths}

\begin{xltabular}{\linewidth}{@{}>{\hsize=0.55\hsize}XX>{\hsize=1.39\hsize}X@{}}
\caption{The path through Metal and the native path below it, compared.}\label{tab:paths-compared}\\
\toprule
& through Metal & below Metal \\
\midrule
\endfirsthead
\tablecontinued{3}
\toprule
& through Metal & below Metal \\
\midrule
\endhead
models & InternLM2.5-1.8B, Qwen3-0.6B, Qwen3-8B & MiniLM-L6, Qwen2.5-0.5B \\*
per token & one command buffer & one Submit per stage (916 for a Qwen decode token) \\*
pipeline configuration & Metal derives it from the image's metadata & the runner writes it, from measured classes \\*
speed & ahead of mlx-lm for some workloads (showcase, section 2; MM~\mm{25.190}) & slower than matched Metal blocks and MPS; no advantage either way
against the full-model Metal control \\*
established & correctness on hardware; performance against mlx-lm & complete execution with no Metal; correctness against the original
checkpoint \\*
not established & why Apple's code for the same kernels is faster & a speedup from bypassing Metal; behaviour across OS updates;
reproduction outside this campaign; a result only this path can produce
(MM~\mm{25.210.9}) \\*
\bottomrule
\end{xltabular}

Performance work is done on the Metal path, which the native path does not replace.
